\chapter{总结与展望}

\section{工作总结}

本文针对水下动态场景的新视角合成问题，将水下物理建模与动态高斯溅射框架相结合，提出了统一的重建系统。主要工作包括：

（1）将SeaSplat的水下物理建模模块（衰减网络、后向散射网络及损失函数）集成至4DGaussians框架，实现对水下光学退化与动态场景变化的联合处理。

（2）实现基于Depth Anything V2的单目深度监督机制，通过逆深度归一化与alpha加权损失提升几何一致性，抑制浮游点云伪影。

（3）设计多阶段渐进式训练策略，通过独立水下优化器、后向散射与前向衰减参数预热、后向散射与前向衰减参数与高斯参数交替更新缓解多模块协同训练的梯度冲突，提升训练稳定性。

实验表明，本文方法在四个动态水下场景上平均PSNR达30.54 dB，较UDR-GS提升1.6\%，较4DGS提升3.1\%。其中Robot、Coral和Streaks场景取得最优PSNR，Fish场景仍受高速非刚性运动影响，说明当前物理建模与动态形变的联合优化仍存在改进空间。

\section{未来工作}

未来可从以下几个方面继续改进：

（1）\textbf{自适应训练调度}。当前训练的阶段切换时机和交替更新频率是人为设定的固定值，换一个场景可能需要重新调整。一个可行的思路是根据训练过程中的梯度变化或验证集上的重建误差来自动决定何时切换阶段、以及水下模块与高斯模块的更新频率比例，从而减少手动调参的工作量。

（2）\textbf{时序信息的利用}。本文目前对每帧图像独立进行处理，没有显式利用相邻帧之间的运动关联。Fish场景的结果表明，当鱼体运动较快时，水下物理模块与变形模块之间会产生一定的优化冲突。一个可能的改进方向是加入光流等时序约束，让相邻帧之间共享一部分运动信息，帮助系统在快速运动区域更稳定地估计水体参数。

（3）\textbf{更精细的物理模型}。当前的水下成像模型假设整个场景范围内的水体光学参数是均匀的，没有考虑实际水体可能存在分层或局部浑浊的情况。后续可以将逐通道的固定衰减系数改为随空间位置变化的形式，使其能够适应更复杂的水下环境。此外，如果能够获取偏振或光谱信息，也可以将现有的RGB三通道模型扩展为更丰富的多光谱模型。

（4）\textbf{水下场景专用的深度估计}。本文使用的Depth Anything V2虽然在通用场景上表现不错，但其训练数据以陆地场景为主，对水下散射模糊、低对比度区域的深度估计效果还有提升空间。后续可以收集水下双目或结构光数据对其进行微调，或者尝试将深度估计网络和高斯渲染管线一起训练，让深度估计随场景重建的改善而更新。

（5）\textbf{推理速度优化}。当前系统的推理时间主要由变形网络前向计算、水下网络逐像素处理以及高斯光栅化三步构成，还达不到水下机器人实时导航的要求。可以从几个方面入手：剪除对渲染贡献不大的冗余高斯点、压缩变形MLP的结构、将水下物理模型简化为更轻量的计算形式。如果条件允许，还可以尝试将整个渲染流程部署到NVIDIA Jetson等嵌入式计算平台上，利用TensorRT等工具进行推理加速。

（6）\textbf{更多场景下的验证与快速适配}。目前仅在四个动态水下场景和四个静态水下场景上进行了测试，后续需要在更广泛的水下环境中验证方法的稳定性，例如高浑浊的近岸水域、深海弱光环境、不同水色类型的水域等。此外，还可以尝试让系统在遇到新场景时，仅利用场景自身的重建误差对水下物理模型做少量迭代微调，使其快速适应新场景的水体特性，而不需要重新进行完整训练。
